\scanpage{323}
\mspage{273}

$\deleted{\unsure{b}}2^\circ)$ \deleted{Si dans $(\mathcal M,W)$, \unsure{on a}} \inserted{Supposons que $\mathcal M$ ait un objet final $e$. Alors} l'implication

\[
b\mathbin{\&}c\Longrightarrow a
\]

pour \uline{tout} hom $f:X\longrightarrow Y$ de carrés $W$-cocartésiens équivaut à l'axiome suivant sur $W$ :

\unsure{(}\uline{Axiome des carrés \inserted{$W$-}cocartésiens conservatifs} : Si

\diagram{\(\begin{tikzcd}[ampersand replacement=\&,column sep=small,row sep=small]
{}\&Z_1\arrow[dr,dash,"{j_2}"]\&{}\\Z_0\arrow[ur,"{i_1}"]\arrow[dr,"{i_2}"']\&{}\&Z_3\\{}\&Z_2\arrow[ur,"{j_1}"']\&{}
\end{tikzcd}\)}

% slanted
\margin{Cf argument page 78 qui montre que l'implication $b\mathbin{\&}c\Longrightarrow a$ entraîne l'axiome conservatif}

$W$-cocartésien, alors $j_1\in W\Longrightarrow i_1\in W$ (donc en fait $i_1\in W\Longleftrightarrow j_1\in W$ et $i_2\in W\Longleftrightarrow j_2\in W$).

Cet axiome est équivalent à la conjonction des deux axiomes \unsure{suivant|suivants}

\uline{Cons. 1} Si on a un carré \deleted{\unsure{$W$-}}cocartésien

\diagram{\(\begin{tikzcd}[ampersand replacement=\&,column sep=small,row sep=small]
\deleted{\unsure{X}}S\arrow[r,"f"]\arrow[dr,dash]\&\deleted{\unsure{Z}}T\arrow[dr,"{\deleted{\unsure{\mathrm{ou}}}}"]\&{}\\{}\&e\arrow[r,dash]\&U
\end{tikzcd}\qquad\unsure{(-}\)}

\deleted{(\unsure{où on peut supposer}} \inserted{avec} $f\in\mathrm{Cof}_W$, \deleted{et le carré cartésien alors} alors $U\in\mathrm{Asph}_W\Longrightarrow f\in W$ (donc $U\in\mathrm{Asph}_W\Longleftrightarrow f\in W$)

\uline{Cons. 2} Si on a un carré cartésien,

\diagram{\deleted{avec}\quad\(\deleted{\unsure{S}\longrightarrow}\begin{tikzcd}[ampersand replacement=\&,row sep=small,column sep=small]
{}\&Z_1\arrow[dr,dash]\&{}\\Z_0\arrow[ur,dash,"{i_1}"]\arrow[dr,dash,"{i_2}"']\&{}\&Z_3\\{}\&Z_2\arrow[ur,dash]\&{}
\end{tikzcd}\)}

avec $i_1,i_2\in\mathrm{Cof}_W$, alors si $Z_1,Z_2,Z_3$ sont $W$ asph, $Z_0$ l'est aussi.
% archive sheet number: 323
